\section{Domain Specific Notation}
\label{app:domain-specific-notation}

\begin{definition}[SI unit convention]
\label{def:domain-si-unit-convention}
Unless otherwise stated, physical units are SI: length in meters,
time in seconds, density in $\mathrm{kg}/\mathrm{m}^3$, pressure in Pa,
dynamic viscosity in $\mathrm{Pa}\cdot\mathrm{s}$, and kinematic viscosity in
$\mathrm{m}^2/\mathrm{s}$.
\end{definition}

\begin{convention}[Centimeter--second numerical-record convention]
\label{conv:cm-second-numerical-record}
The SI convention above governs continuum notation and unqualified analytical
parameters. The Julia solver outputs and the
Section~\ref{sec:resolved-3d-comparison} tables instead use the
centimeter--gram--second convention of the numerical realization:
lengths are in $\mathrm{cm}$, time in $\mathrm{s}$, density in
$\mathrm{g}/\mathrm{cm}^3$, velocity in $\mathrm{cm}/\mathrm{s}$, pressure
and stress in $\mathrm{dyn}/\mathrm{cm}^2$, dynamic viscosity in
$\mathrm{g}/(\mathrm{cm}\,\mathrm{s})=\mathrm{dyn}\,\mathrm{s}/\mathrm{cm}^2$,
and kinematic viscosity in $\mathrm{cm}^2/\mathrm{s}$. Thus the default
solver geometry parameters $L$, $R_{\max}$, $z_a$, $z_s$,
$\sigma_a$, and $\kappa$ are centimeter quantities, while the stenosis
localization $\lambda$ in $\exp(-\lambda g^4)$ has units
$\mathrm{cm}^{-4}$. The radius-squared solver state coordinate $a=R^2$ has units
$\mathrm{cm}^2$; when a physical circular area is meant it is
$A_{\mathrm{phys}}=\pi a$, also reported in square centimeters.

Output fields carry their display units in their names. In particular,
\texttt{z\_cm} and \texttt{R0\_cm} are geometry fields, \texttt{A\_cm2} is the
area-coordinate field in $\mathrm{cm}^2$, and \texttt{Q\_cm3\_s} is
the flow-coordinate field paired with $a$ in nominal $\mathrm{cm}^3/\mathrm{s}$.
For the radius-squared solver convention, the corresponding physical
circular-area flow is $Q_{\mathrm{phys}}=\pi Q$, while
\texttt{uavg\_cm\_s} denotes
$Q/a=Q_{\mathrm{phys}}/A_{\mathrm{phys}}$. The
Section~\ref{sec:resolved-3d-comparison} velocity/error columns are in
$\mathrm{cm}/\mathrm{s}$, \texttt{pressure\_dyn\_cm2} and initial-condition
pressure-drop fields are in $\mathrm{dyn}/\mathrm{cm}^2$, and
\texttt{nu\_eff\_cm2\_s} is in $\mathrm{cm}^2/\mathrm{s}$. A pressure supplied
through a Pa-valued command-line option is converted before storage using
$1\,\mathrm{Pa}=10\,\mathrm{dyn}/\mathrm{cm}^2$. Any numerical-output value
used in an SI formula must first be converted; for example
$1\,\mathrm{cm}=10^{-2}\,\mathrm{m}$,
$1\,\mathrm{cm}^2/\mathrm{s}=10^{-4}\,\mathrm{m}^2/\mathrm{s}$,
$1\,\mathrm{g}/(\mathrm{cm}\,\mathrm{s})=10^{-1}\,\mathrm{Pa}\cdot\mathrm{s}$,
and $1\,\mathrm{dyn}/\mathrm{cm}^2=10^{-1}\,\mathrm{Pa}$.
\end{convention}

\begin{convention}[Pressure-output convention]
\label{conv:appendix-pressure-reference-ratio}
Pressure-drop and pressure-ratio quantities are reduced output vocabulary. A
pressure-drop trace requires proximal and distal scalar pressure observations,
and a pressure-ratio output is reported only with a ratio-output specification
that declares proximal and distal scalar pressure observations, a time-averaging
window, a pressure reference, and the denominator admissibility condition below.
The general ratio-output specification is
\begin{flalign*}
    \quad
    \pi_{\mathrm{PR}}
    \defined
    (\mathcal O_{\mathrm{prox}},
    \mathcal O_{\mathrm{dist}},
    \mathcal{T}_{\mathrm{HF}},
    p_{\mathrm{ref}}). &&
\end{flalign*}
Here $\mathcal O_{\mathrm{prox}}$ and $\mathcal O_{\mathrm{dist}}$ are the
declared scalar observation maps applied to the gauge-pressure field
$p_{\mathrm g}$, and $p_{\mathrm{ref}}$ is the fixed scalar pressure reference
added to observed gauge-pressure traces before forming the ratio. In the 1D
axial-tap or cross-sectional
pressure-average specialization, this report may use the shorthand
\begin{flalign*}
    \quad
    \pi_{\mathrm{PR}}
    =
    (z_{\mathrm{prox}},
    z_{\mathrm{dist}},
    \mathcal{T}_{\mathrm{HF}},
    p_{\mathrm{ref}}), &&
\end{flalign*}
with the pressure observation maps understood from the declared tap or
cross-sectional averaging convention.
Given a model gauge-pressure field, the observed pressure traces are
\begin{flalign*}
    \quad
    P_{\mathrm{prox}}(t)
    &\defined
    \mathcal O_{\mathrm{prox}}[p_{\mathrm g}(\cdot,t)], &
    P_{\mathrm{dist}}(t)
    &\defined
    \mathcal O_{\mathrm{dist}}[p_{\mathrm g}(\cdot,t)]. &&
\end{flalign*}
The corresponding pressure-drop trace is
\begin{flalign*}
    \quad
    \Delta p(t)
    \defined
    P_{\mathrm{prox}}(t)-P_{\mathrm{dist}}(t). &&
\end{flalign*}
A pressure ratio is not invariant under common additive pressure shifts. The
reference-consistent ratio traces are therefore defined after observation:
\begin{flalign*}
    \quad
    P_{\mathrm{prox}}^{\mathrm{ratio}}(t)
    &\defined
    p_{\mathrm{ref}}+P_{\mathrm{prox}}(t), &
    P_{\mathrm{dist}}^{\mathrm{ratio}}(t)
    &\defined
    p_{\mathrm{ref}}+P_{\mathrm{dist}}(t). &&
\end{flalign*}
This trace-level convention does not require the observation maps to preserve
constant shifts. The subscript $\mathrm{HF}$ labels a declared high-flow
averaging window in the
study record; it is not a clinical hyperemia model unless an external
measurement comparison is separately specified
\cite[Ch.~5, Sec.~5.5, pp.~72--73]{EscanedDavies2017PhysiologicalAssessment}
\cite[p.~1703]{PijlsEtAl1996FFRFunctionalSeverity}.
The ratio-output specification is admissible for a computed pressure trace when
$\mathcal{T}_{\mathrm{HF}}\subset I_T$ is Lebesgue-measurable with
$0<|\mathcal{T}_{\mathrm{HF}}|<\infty$, the proximal and distal
ratio-pressure traces are integrable on that window, and
\begin{flalign*}
    \quad
    \langle P_{\mathrm{prox}}^{\mathrm{ratio}}(t)
    \rangle_{\mathcal{T}_{\mathrm{HF}}}
    \ne 0. &&
\end{flalign*}
For an admissible $\pi_{\mathrm{PR}}$, the corresponding reduced
pressure-ratio output is
\begin{flalign*}
    \quad \mathrm{PR}_{1D}(\pi_{\mathrm{PR}}) \defined
    \frac{\langle P_{\mathrm{dist}}^{\mathrm{ratio}}(t)\rangle_{\mathcal{T}_{\mathrm{HF}}}}
         {\langle P_{\mathrm{prox}}^{\mathrm{ratio}}(t)\rangle_{\mathcal{T}_{\mathrm{HF}}}}. &&
\end{flalign*}
For an integrable scalar trace,
\begin{flalign*}
    \quad \langle f\rangle_{\mathcal{T}_{\mathrm{HF}}} \defined \frac{1}{|\mathcal{T}_{\mathrm{HF}}|} \int_{\mathcal{T}_{\mathrm{HF}}} f(t)\,dt. &&
\end{flalign*}
This bracket denotes time averaging, not the inner product notation in
Appendix~\ref{app:mathematical-notation}.
An FFR-like shorthand refers to this reduced ratio only for records carrying an
admissible $\pi_{\mathrm{PR}}$; the shorthand does not add a clinical
measurement model
\cite[Ch.~5, Sec.~5.5, pp.~72--73]{EscanedDavies2017PhysiologicalAssessment}
\cite[p.~1703]{PijlsEtAl1996FFRFunctionalSeverity}. If $\pi_{\mathrm{PR}}$ is
absent or inadmissible, the pressure ratio is outside the reported output space.
\end{convention}

\begin{convention}[Shear-output convention]
\label{conv:appendix-shear-diagnostic}
Shear terminology is model-tier dependent. A resolved wall shear stress output
requires a resolved stress tensor, wall normal, wall location, output operator,
and, for a signed component, a chosen tangent direction
\cite[Abstract; Sec.~2, pp.~8--10]{JohnEtAl2017WallShearStressVectorForm}. A
reduced shear proxy such as $\tau_{w,z}^{1D}$ requires the model tier, wall
closure, rheology descriptor, velocity-profile or friction closure, evaluation
location or spatial averaging operator, and sign or time-window convention.
Reduced shear proxies are closure diagnostics; they are not interchangeable with
resolved traction-based wall shear stress unless a separate comparison map is
declared.
\end{convention}


\begin{minipage}[t]{\textwidth}
    \begin{tabular}{@{}p{0.22\linewidth} p{0.50\linewidth} p{0.18\linewidth}@{}}
        \toprule
        \textbf{Symbol} & \textbf{Meaning} & \textbf{Units} \\
        \midrule
        $R$ & vessel radius; vessel diameter is $d_v := 2R$ & $\mathrm{m}$ \\
        $d_v$ & vessel diameter & $\mathrm{m}$ \\
        $\eta$ & dynamic viscosity & $\mathrm{Pa}\cdot\mathrm{s}$ \\
        $\nu := \eta/\rho$ & kinematic viscosity & $\mathrm{m}^2/\mathrm{s}$ \\
        $\eta_{\mathrm{eff}}(\dot{\gamma})$
            & effective dynamic viscosity for generalized Newtonian models
            & $\mathrm{Pa}\cdot\mathrm{s}$ \\
        $\nu_{\mathrm{eff}}(\dot{\gamma})
        :=\eta_{\mathrm{eff}}(\dot{\gamma})/\rho$
            & effective kinematic viscosity after density scaling
            & $\mathrm{m}^2/\mathrm{s}$ \\
        $\tau$ & shear stress & $\mathrm{Pa}$ \\
        $\dot{\gamma}$ & shear rate & $\mathrm{s}^{-1}$ \\
        $\rho$ & density field & $\mathrm{kg}/\mathrm{m}^3$ \\
        $p$ & pressure field & $\mathrm{Pa}$ \\
        $\vu$ & velocity field & $\mathrm{m}/\mathrm{s}$ \\
        $D(\vu)$ & rate-of-deformation tensor,
        $D(\vu) := \frac{1}{2}\left(\nabla \vu + \nabla \vu^\top\right)$
        & $\mathrm{s}^{-1}$ \\
        $W_0:=d_v\sqrt{\frac{\omega\rho}{\eta}}$ & Womersley number, diameter-based convention & $-$ \\
        $\mathrm{Re}$ & Reynolds number & $-$ \\
        $\mathrm{Pe}$ & Péclet number & $-$ \\
        $c_{\mathrm{conc}}$ & concentration field when a transport scalar is discussed & $-$ \\
        $D_m$ & molecular diffusion coefficient & $\mathrm{m}^2/\mathrm{s}$ \\
        $t$ & time & $\mathrm{s}$ \\
        $T$ or $T_{\mathrm{final}}$
            & terminal time; use $T_{\mathrm{final}}$ where it could be
            confused with stress tensors
            & $\mathrm{s}$ \\
        $\omega$ & angular frequency & $\mathrm{rad}/\mathrm{s}$ \\
        $(r,\theta,z)$
            & cylindrical coordinates used only when 3D coordinate forms are
            written
            & $\mathrm{m},-,\mathrm{m}$ \\
        $(\vct{e}_r,\vct{e}_\theta,\vct{e}_z)$ & cylindrical orthonormal basis vectors & $-$ \\
        $\tau_{ij}$ & cylindrical viscous-stress component & $\mathrm{Pa}$ \\
        $\vct{f_b}$ & body force per unit mass & $\mathrm{m}/\mathrm{s}^2$ \\
        $\rho\vct{f_b}$ & body force per unit volume & $\mathrm{N}/\mathrm{m}^3$ \\
        \bottomrule
\end{tabular}
\end{minipage}

\par\smallskip\noindent
The symbols $\mathrm{Pe}$, $c_{\mathrm{conc}}$, and $D_m$ are reserved for a
future passive-scalar transport extension, for example
\begin{flalign*}
    \quad
    \pd_t c_{\mathrm{conc}}+\vu\cdot\nabla c_{\mathrm{conc}}
    =
    D_m\Delta c_{\mathrm{conc}},
    \qquad
    \mathrm{Pe}=\frac{UL}{D_m},
    &&
\end{flalign*}
where $\mathrm{Pe}$ compares advective and diffusive scalar transport.

\paragraph*{Reduced-model and forward-map notation.}
\par\smallskip\noindent

\begingroup
\small
\begin{longtable}{@{}>{\raggedright\arraybackslash}p{0.22\linewidth}
        >{\raggedright\arraybackslash}p{0.54\linewidth}
        >{\raggedright\arraybackslash}p{0.14\linewidth}@{}}
        \toprule
        \textbf{Symbol} & \textbf{Meaning} & \textbf{Units} \\
        \midrule
        \endfirsthead
        \toprule
        \textbf{Symbol} & \textbf{Meaning} & \textbf{Units} \\
        \midrule
        \endhead
        $A(z,t)$ & cross-sectional lumen area in the 1D model & $\mathrm{m}^2$ \\

        $Q(z,t)$ & volumetric flow rate through a cross-section & $\mathrm{m}^3/\mathrm{s}$ \\

        $\bar u(z,t):=Q/A$ & section-mean axial velocity & $\mathrm{m}/\mathrm{s}$ \\

        $I_z=(0,L)$ & axial single-vessel interval for a 1D formulation & $\mathrm{m}$ \\

        $\mathbf U=(A,Q)^\top$ & conservative area-flow state vector & $(\mathrm{m}^2,\mathrm{m}^3/\mathrm{s})$ \\

        $\mathcal P$
            & cross-sectional velocity-profile descriptor used by the 1D
            closure; admitted values include flat, parabolic, and
            power-family profiles
            & mixed \\

        $\alpha_{\mathcal P}$
            & momentum-flux correction induced by the declared profile
            $\mathcal P$; the parabolic baseline uses
            $\alpha_{\mathcal P}=4/3$
            & $-$ \\

        $\gamma$
            & power-profile exponent in
            $u(r)=((\gamma+2)/\gamma)\bar u(1-(r/R)^\gamma)$
            & $-$ \\

        $g_{\mathcal P}$
            & profile shear-rate factor used by the 1D rheology proxy; for
            power profiles $g_{\mathcal P}=\gamma+2$, while flat profiles
            assign this value explicitly
            & $-$ \\

        $R(A(z,t))=\sqrt{A(z,t)/\pi}$
            & current reduced lumen radius derived from the 1D area
            & $\mathrm{m}$ \\

        $R_0(z)$ & reference stenosed radius used by the wall law & $\mathrm{m}$ \\

        $R_{\mathrm{base}}(z)$ & nominal healthy baseline radius before applying $s(z)$ & $\mathrm{m}$ \\

        $\mathcal R$
            & rheology descriptor in the 1D model; distinct from radius symbols
            $R$, $R_0$, and $R_{\mathrm{base}}$; admitted values are
            \texttt{newtonian}, \texttt{carreau},
            \texttt{carreau-yasuda}, \texttt{casson}, and \texttt{power-law}
            & mixed \\

        $\mathcal W$
            & wall or FSI closure descriptor; the selected 1D report model uses
            $\mathcal W_{\mathrm{elastic1D}}$, while fixed-wall,
            multiscale-coupled, and resolved-FSI models require separate
            declarations
            & mixed \\

        $s(z)$ & dimensionless stenosis profile, with $0\leq s(z)<1$ & $-$ \\

        $A_0(z)=\pi R_0(z)^2$ & wall-law reference area after applying the stenosis profile & $\mathrm{m}^2$ \\

        $s_{\max}$ & maximum fractional radius-reduction parameter in $s(z)$ & $-$ \\

        $g(z)$ & smooth asymmetric stenosis-kernel coordinate & $\mathrm{m}$ \\

        $\psi(z)=\exp(-\lambda g(z)^4)$
            & smooth rapidly localized stenosis kernel used by the baseline
            profile $s(z)=s_{\max}\psi(z)$
            & $-$ \\

        $z_a$ & center of the Gaussian asymmetry term in $g(z)$ & $\mathrm{m}$ \\

        $z_s$ & axial shift parameter in $g(z)$ & $\mathrm{m}$ \\

        $\sigma_a$ & width parameter for the Gaussian asymmetry term in $g(z)$ & $\mathrm{m}$ \\

        $\kappa$ & amplitude of the Gaussian asymmetry term in $g(z)$ & $\mathrm{m}$ \\

        $\lambda$ & localization strength in $\psi(z)=\exp(-\lambda g(z)^4)$ & $\mathrm{m}^{-4}$ \\

        $\beta(z)$ & wall-stiffness parameter in
        $p_{\mathrm{g}}=p_{\mathrm{ext}}+\beta A_0^{-1}(\sqrt A-\sqrt{A_0})$ & $\mathrm{Pa}\cdot\mathrm{m}$ \\

        $p_{\mathrm{ext}}(z,t)$
            & external-pressure value in the wall-law gauge convention;
            baseline value $p_{\mathrm{ext}}\equiv0$ unless declared otherwise
            & $\mathrm{Pa}$ \\

        $\Psi(A,z,t)$ or $\Psi_{\mathcal W}(A,z,t)$
            & elastic flux potential satisfying
            $\pd_A \Psi_{\mathcal W}
            =(A/\rho) \pd_A p_{\mathrm{g}}^{\mathcal W}$ and normalized by
            $\Psi_{\mathcal W}(A_0(z),z,t)=0$ when that normalization is
            declared
            & $\mathrm{m}^4/\mathrm{s}^2$ \\

        $\mathbf F$ & area-flow flux map; not a residual or objective functional & problem-dependent \\

        $\mathbf S$ & source-term map in the 1D balance-law template & problem-dependent \\

        $S_{\mathrm{geom}}^{\mathcal W}$
            & fixed-area wall/geometry source associated with the selected
            reduced wall closure $\mathcal W$
            & $\mathrm{m}^3/\mathrm{s}^2$ \\

        $S_f^{\mathcal R}$
            & wall-friction source associated with the selected rheology or
            effective-viscosity closure $\mathcal R$
            & $\mathrm{m}^3/\mathrm{s}^2$ \\

        $\widehat{\mathbf F}_{i+1/2}^n$
            & numerical flux at a finite-volume cell interface and time level
            & problem-dependent \\

        $\widehat{\mathbf S}_i^n$
            & cell source quadrature used by a computed finite-volume scheme
            & problem-dependent \\

        $\Delta z_i,\Delta t^n$
            & finite-volume cell length and timestep in a computed solve
            & $\mathrm{m},\mathrm{s}$ \\

        $\dot\gamma_{1D}$
            & computed shear-rate proxy
            $g_{\mathcal P} |Q|/(a\sqrt a)$, where $a=R^2$
            & $\mathrm{s}^{-1}$ \\

        $\nu_{\mathrm{eff}}^{\mathcal R}$
            & closure-indexed effective kinematic viscosity used in the
            reduced viscous/source and pressure-correction terms
            & $\mathrm{m}^2/\mathrm{s}$ \\

        $C_f$ or $C_f(z,A)$ & wall-friction coefficient in the 1D source
        $-C_fQ/A$; constant-viscosity value $C_f=8\pi\eta/\rho$ & $\mathrm{m}^2/\mathrm{s}$ \\

        $\eta_{\mathrm{fric}}(z,t)$
            & dynamic viscosity used by the reduced wall-friction and
            signed-shear-proxy formula; $\eta_{\mathrm{fric}}=\eta$ for
            constant-viscosity realizations and
            $\eta_{\mathrm{fric}}=\eta_{\mathrm{eff}}^{\mathcal R}
            (\dot\gamma_{1D})$ for admitted non-Newtonian descriptor
            realizations
            & $\mathrm{Pa}\cdot\mathrm{s}$ \\

        $c_{\mathrm{wave}}$ or $c$ & 1D pulse-wave speed,
        $c^2=(A/\rho) \pd_A p_{\mathrm{g}}$ when no concentration field is in scope & $\mathrm{m}/\mathrm{s}$ \\

        $e\in\mathcal E$
            & expert/closure record for a closure-indexed 1D forward-map
            family; suppress $e$ only for single-expert baseline records
            & mixed \\

        $\mathcal G_{1D,h}^{r_h,e}$
            & computed fixed-grid 1D map indexed by numerical record $r_h$
            and expert record $e$
            & mixed \\
        \bottomrule
\end{longtable}
\endgroup


\begin{minipage}[t]{\textwidth}
    \footnotesize
    \begin{tabular}{@{}>{\raggedright\arraybackslash}p{0.22\linewidth}
        >{\raggedright\arraybackslash}p{0.54\linewidth}
        >{\raggedright\arraybackslash}p{0.14\linewidth}@{}}
        \toprule
        \textbf{Symbol} & \textbf{Meaning} & \textbf{Units} \\
        \midrule

        $p_{3D}$ & continuum pressure in the 3D/NS setting & $\mathrm{Pa}$ \\

        $p_{3D,\mathrm{g}}$
            & continuum pressure after subtracting the wall-law gauge reference
            & $\mathrm{Pa}$ \\

        $p_{\mathrm{g}}$ & 1D gauge pressure produced by the wall law & $\mathrm{Pa}$ \\

        $p_{\mathrm{ratio}}:=p_{\mathrm{ref}}+p_{\mathrm{g}}$
            & field-level reference-consistent pressure shorthand; ratio
            traces are defined after scalar observation by
            Convention~\ref{conv:appendix-pressure-reference-ratio}
            & $\mathrm{Pa}$ \\

        $p_{\mathrm{ref}}$
            & fixed pressure reference recorded by the ratio-output
            specification
            & $\mathrm{Pa}$ \\

        $\mathcal O_{\mathrm{prox}},\mathcal O_{\mathrm{dist}}$
            & proximal and distal scalar pressure observation maps declared
            for pressure-drop or pressure-ratio outputs
            & input/output $\mathrm{Pa}$ \\

        $P_{\mathrm{prox}},P_{\mathrm{dist}}$
            & observed proximal and distal gauge-pressure traces after applying
            the declared pressure observation maps
            & $\mathrm{Pa}$ \\

        $P_{\mathrm{prox}}^{\mathrm{ratio}},
        P_{\mathrm{dist}}^{\mathrm{ratio}}$
            & observed proximal and distal pressure traces after adding
            $p_{\mathrm{ref}}$
            & $\mathrm{Pa}$ \\

        $\pi_{\Delta p}=(z_{\mathrm{prox}},z_{\mathrm{dist}})$
            & 1D shorthand pressure-drop tap specification; the observation
            maps define the actual scalar traces
            & $\mathrm{m},\mathrm{m}$ \\

        $\pi_{\mathrm{PR}}$
            & ratio-output specification
            $(\mathcal O_{\mathrm{prox}},
            \mathcal O_{\mathrm{dist}},\mathcal{T}_{\mathrm{HF}},
            p_{\mathrm{ref}})$; axial taps may be specified by shorthand
            & mixed \\

        $\langle f\rangle_{\mathcal{T}_{\mathrm{HF}}}$
            & time average of an integrable trace over the declared high-flow
            window; not the Appendix~\ref{app:mathematical-notation} inner
            product
            & same as $f$ \\

        $\langle P_{\mathrm{prox}}^{\mathrm{ratio}}\rangle_{\mathcal{T}_{\mathrm{HF}}}\ne0$
            & denominator admissibility condition required before
            $\mathrm{PR}_{1D}$ is reported
            & $-$ \\

        $p_{\mathrm{out},\mathrm{g}}$
            & prescribed outlet gauge pressure in the same reference frame as
            $p_{\mathrm{g}}$
            & $\mathrm{Pa}$ \\

        $\Delta p(t)$
            & proximal-minus-distal scalar pressure trace after the declared
            observations
            & $\mathrm{Pa}$ \\

        $\overline{\Delta p}_{\mathcal{T}_{\mathrm{HF}}}$
            & scalar pressure-drop output averaged over the declared high-flow
            window $\mathcal{T}_{\mathrm{HF}}$
            & $\mathrm{Pa}$ \\

        FFR-like shorthand
            & optional shorthand for $\mathrm{PR}_{1D}$ only when
            Convention~\ref{conv:appendix-pressure-reference-ratio}
            and admissibility requirements are met
            & $-$ \\

        $\mathrm{PR}_{1D}(\pi_{\mathrm{PR}})$
            & reduced pressure-ratio output under the declared observation
            maps, pressure reference, and averaging window
            & $-$ \\

        $\tau_w$ or $|\tau_w|$
            & resolved traction-based wall-shear-stress output after declaring
            the stress tensor, wall normal, location, and output operator
            & $\mathrm{Pa}$ \\

        $\tau_{w,z}^{1D}$
            & signed 1D axial shear proxy after declaring the model tier,
            closures, location, and sign or window convention
            & $\mathrm{Pa}$ \\

        \texttt{shear\_rate\_1\_s}
            & CSV output field for the computed shear-rate proxy
            $\dot\gamma_{1D}$
            & $\mathrm{s}^{-1}$ \\

        \texttt{nu\_eff\_cm2\_s}
            & CSV output field for the effective kinematic viscosity used by
            the declared rheology closure
            & $\mathrm{cm}^2/\mathrm{s}$ \\

        \texttt{rheology}
            & CSV and summary output field naming the declared rheology
            descriptor
            & $-$ \\
        \bottomrule
    \end{tabular}
\end{minipage}

\begin{remark}[Viscosity notation]
\label{rmk:viscosity-notation}
We write $\eta$ for dynamic viscosity and
\begin{flalign*}
    \quad \nu := \frac{\eta}{\rho} &&
\end{flalign*}
for kinematic viscosity. The distinction is dimensional: this report follows the
convention that the
coefficient in the Cauchy stress tensor is a dynamic viscosity: $D(\vu)$ has
units of inverse time, so $\eta D(\vu)$ has pressure units. For generalized
Newtonian models, the scalar viscosity function multiplying $2D(\vu)$ in the
extra stress is likewise an effective dynamic viscosity, denoted
\begin{flalign*}
    \quad \eta_{\mathrm{eff}}(\dot{\gamma}), \qquad \dot{\gamma}\defined\sqrt{2\vct{D}(\vu):\vct{D}(\vu)}. &&
\end{flalign*}
Here ``this quantity'' means the stress-level viscosity coefficient, not the
divided-by-density coefficient. Several sources write the generalized viscosity
function as $\mu(|D|^2)$, while related hemodynamics sources also use $\mu$
for viscosity coefficients in reduced and continuum settings
\cite[Sec.~3.1, p.~35; Sec.~3.6, p.~215]{KopplHelmig2023DimensionReducedModeling}
\cite[Sec.~1.1, pp.~6--7; Eqs.~(1.2)--(1.4); Fig.~4.4, p.~18]{JohnEtAl2017WallShearStressVectorForm}.
To keep this report unambiguous, $\nu$ is always kinematic viscosity, while
$\eta$ and $\eta_{\mathrm{eff}}$ are dynamic viscosities.
\end{remark}
